\documentclass[tikz,border=6pt]{standalone}
\usepackage{amsmath}
\usetikzlibrary{arrows.meta,calc,angles,quotes}
\begin{document}
\begin{tikzpicture}[>=Latex, font=\sffamily]
  \node[anchor=west, text=blue!60!black, font=\bfseries\Large] at (-0.3,4.3) {Torque, momento angular y momento de inercia};
  \draw[rounded corners=8pt, fill=blue!3, draw=blue!20] (-0.4,-0.4) rectangle (10.5,3.8);
  \fill (1.3,1.4) circle (2pt) node[below=4pt]{eje};
  \coordinate (O) at (1.3,1.4); \coordinate (P) at (4.8,2.1);
  \draw[->, very thick, blue] (O) -- (P) node[midway, above] {$\vec r$};
  \draw[->, very thick, red] (P) -- ++(0,1.8) node[right] {$\vec F$};
  \draw pic[draw=orange, thick, "$\theta$", angle radius=0.8cm] {angle = P--O--$(P)+(0,1)$};
  \node at (2.2,0.1) {$\vec\tau=\vec r\times\vec F$};
  \draw[blue, thick] (7.7,1.5) circle (1.0);
  \draw[->, very thick, purple] (7.7,0.3) -- (7.7,3.1) node[right] {$\vec\omega$};
  \draw[->, very thick, green!60!black] (8.5,2.1) arc (25:260:0.8);
  \node at (7.7,0.0) {$L=I\omega$};
\end{tikzpicture}
\end{document}
